\section{About the Topic}

\subsection{Key Terms}

\textbf{Protectionism} – The use of government policies such as tariffs, subsidies, import quotas or other restrictions placed on the imports of foreign competitors to shield domestic industries. While protectionism may serve short-term economic or political goals, it is broadly in tension with the WTO’s foundational commitment to open and non-discriminatory trade.

\textbf{Most-Favoured  Nation  (MFN)}  –  A  core non-discrimination principle of the WTO (see GATT Article I) requiring that any trade advantage granted by one member to another must be extended equally to all WTO members. MFN treatment ensures no trading partner is disadvantaged relative to any other without a legal justification.

\textbf{National  Treatment}  –  A  complementary non-discrimination principle of the WTO (see GATT Article III) mandating that once imported goods have entered a market, they must be treated no less favourably than equivalent domestically produced products. National treatment prevents governments from using domestic regulation as a covert form of protectionism.

\textbf{GATT Article XXI} — Article XXI “The Essential Security Exception" permits member states to deviate from their trade obligations when they consider it necessary for the protection of their “essential security interests”. It remains the most contested provision in contemporary WTO law, specifically the scope of this exception and whether it is subject to independent review.

\textbf{Export Controls} – Government-imposed restrictions on the sale, transfer, or re-export of specific goods, technologies or software to foreign buyers, typically justified on national security or foreign policy grounds. Export controls have expanded dramatically since 2022, particularly in the semiconductor and advanced technology industries.

\textbf{Dual-Use Goods} – Products, materials, software, or technologies that have both legitimate civilian applications and potential military or strategic uses. The classification and regulation of dual-use goods is a central challenge in balancing open trade with national security imperatives.

\textbf{Dispute Settlement Body (DSB)} – the WTO organ responsible for administering the rules and procedures governing the resolution of trade disputes between member states. The DSB has the authority to establish panels, adopt panel and Appellate Body reports, and authorise retaliatory measures in the event of non-compliance.

\textbf{Appellate Body} – A standing seven-member tribunal that hears appeals from panel rulings issued by the DSB. Since December 2019, the Appellate Body has been non-functional due to the United States blocking the appointment of new members, creating a systemic enforcement gap within the WTO dispute settlement system.

\textbf{De-risking vs. Decoupling} – Two distinct but related concepts describing how states are restructuring their economic relationships with geopolitical rivals. Decoupling implies a comprehensive severance of economic ties; de-risking refers to a more targeted reduction of strategic dependencies while maintaining broader trade relationships. The distinction has become central to EU-China and US-China policy debates.

\textbf{Industrial Policy} – State intervention designed to develop, support, or protect strategically important domestic industries, often through subsidies, tax incentives, or public investment. Recent examples include the United States' CHIPS and Science Act (2022) and the European Union's Green Deal Industrial Plan, both of which have raised questions about WTO subsidy rules.

\textbf{Economic Coercion} – The deliberate use of trade or economic measures by one state to pressure another into changing its political behaviour or policy positions, without resorting to military force. Economic coercion occupies a grey zone in WTO law, as it may involve measures that technically comply with trade rules while being politically motivated.

\textbf{Friendshoring / Nearshoring} – The restructuring of global supply chains to prioritise trade and production within networks of politically aligned or geographically proximate partners. Friendshoring reflects a growing preference for supply chain security over pure economic efficiency and poses significant questions for the WTO's MFN obligations.

\subsection{Introduction to the Topic}

\subsubsection{A System Built on an Unresolved Tension}

The world trading system was built on a foundational premise: that open, rules-based trade promotes both economic prosperity and peace between nations. Since the establishment of GATT in 1947 and the creation of the WTO in 1995, this premise has underpinned a remarkable expansion of global commerce. Yet the system has always contained a deliberate tension at its core. Alongside its commitments to non-discrimination and market openness, GATT Article XXI permitted member states to deviate from their trade obligations to protect their essential security interests.\footnote{ARTICLE XXI Security Exceptions. (n.d.). In Analytical index of the GATT (pp. 599–610). \url{https://www.wto.org/english/res_e/booksp_e/gatt_ai_e/art21_e.pdf}} For decades, this provision remained largely unused. Today, it sits at the

Industrial Policy and the Return of Protectionism Geopolitical rivalry is driving a wider adoption of industrial policy as a mainstream instrument of statecraft. The United States' CHIPS and Science Act, the European Union's Green Deal Industrial Plan, and China's Made in China 2025 strategy all involve substantial state subsidies directed at strategic industries, raising serious questions under WTO subsidy rules.\footnote{Maruyama, W., \& Wolff, A. Wm. (2023). Saving the WTO from the national security exception. Peterson Institute for International Economics.} Simultaneously, the restructuring of supply chains along political lines through friendshoring and nearshoring is fragmenting the integrated global economy that the WTO was designed to sustain. These trends put developing nations in a particularly vulnerable position as they risk being excluded from the supply chains and centre of the most consequential crisis in the WTO's history.\footnote{Maruyama, W., \& Wolff, A. Wm. (2023). Saving the WTO from the national security exception. Peterson Institute for International Economics.}

investment flows that major powers are redirecting among themselves.\footnote{Long, T. (2026, April 1). Decoupling risks: How semiconductor export controls could harm US chipmakers and innovation. ITIF. \url{https://itif.org/publications/2025/11/10/decoupling-risks-semiconductor-export-controls-har} m-us-chipmakers-innovation/}

\subsubsection{The Abuse of the National Security Exception}

From GATT's founding in 1947 until 1994, Article XXI was formally invoked in only one dispute. Since 1995, there have been fourteen requests for panel establishment on national security grounds, the overwhelming majority of which were filed after 2019, marking a movement towards protectionism.\footnote{Mavroidis, P. C. (n.d.). Litigating national security in the WTO era. Columbia Law School Scholarship Archive. \url{https://scholarship.law.columbia.edu/faculty_scholarship/4656/}} Article XXI grants member states broad discretion to determine what constitutes an essential security interest, creating a loophole that, if left unchecked, risks  hollowing  out  the  WTO's  entire non-discrimination framework. Allowing members to retain complete autonomy over the use of exceptions could lead to a surge in protectionist measures masquerading as security policies, fundamentally undermining the very predictability and openness that the WTO was designed to guarantee.\footnote{The "National Security Exception" and the World Trade Organization. (n.d.). Congress.gov | Library of Congress. \url{https://www.congress.gov/crs-product/LSB10223}}

\subsubsection{The Collapse of the Dispute Settlement System}

The WTO's ability to rule on disputes has been severely compromised. Since December 2019, the Appellate Body has been non-functional, leaving appeals filed by member states unresolved and stripping panel rulings of their binding legal force (Mavroidis, n.d.). The United States, whose blocking of Appellate Body appointments caused the crisis in the first place, went further still, formally declaring in January 2023 that it had been a mistake to allow national security measures to be brought before the WTO at all.\footnote{Lester, S. (2025, September 9). The U.S. proposal on how to handle disputes on essential security measures. International Economic Law and Policy Blog. \url{https://ielp.worldtradelaw.net/2024/12/the-us-proposal-on-how-to-handle-disputes-on-ess} ential-security-measures/} This position, if accepted, would transform the WTO from a binding legal system into one in which the most powerful members can opt out of their obligations at will.

\subsubsection{Semiconductors and Economic Decoupling}

The tension between national security and free trade is nowhere more acute than in the technology sector. Advanced semiconductors are fundamental to artificial intelligence, defence systems and telecommunications infrastructure. Policymakers now consider semiconductor supply chains to be a matter of national survival rather than commercial competition. As of 2022, the United States has progressively tightened chip export controls, with the Trump administration blacklisting dozens of additional Chinese entities in March 2025.\footnote{Reinsch, W. A., \& Caporal, J. (2025, November 17). The WTO's first ruling on national security: What does it mean for the United States? CSIS. \url{https://www.csis.org/analysis/wtos-first-ruling-national-security-what-does-it-mean-united-states}} In response, China has introduced its own export controls on dual-use materials and established a \$47.5 billion domestic chip investment fund, accelerating the process of technological decoupling, which sits uneasily within the WTO's foundational principles of non-discrimination and market openness.\footnote{Long, T. (2026, April 1). Decoupling risks: How semiconductor export controls could harm US chipmakers and innovation. ITIF. \url{https://itif.org/publications/2025/11/10/decoupling-risks-semiconductor-export-controls-har} m-us-chipmakers-innovation/}

At the heart of this committee's debate lies a fundamental conflict between two visions of global trade governance: on one side, a rules-based multilateral order where no state can unilaterally define its own obligations. On the other hand, there is a sovereignty-first security logic, in which states assert the right to act on national security grounds, regardless of their commitments to the WTO. The central challenge delegates must address is striking a workable balance between these two logics, one that takes genuine security concerns seriously without allowing security to become an all-purpose justification for protectionism.

\subsection{History of the Topic}

The origins of this tension of national security and free trade in an era of geopolitical tension go back to the founding of the multilateral trading system, GATT, itself. GATT’s roots lie in the 1944 Bretton Woods Conference, which sought to stabilise post-war economic order after the catastrophic 1930s, a decade defined by competitive devaluations, protectionist spirals and the economic nationalism that contributed to the Second World War. When GATT entered into force in 1948, Article XXI was included as a safety net, recognising that states would never accept trade commitments that left them unable to respond to dire security threats.\footnote{ARTICLE XXI Security Exceptions. (n.d.). In Analytical index of the GATT (pp. 599–610). \url{https://www.wto.org/english/res_e/booksp_e/gatt_ai_e/art21_e.pdf}} It was considered a last resort, not a routine policy tool. By the time the WTO replaced GATT in 1995, 125 nations had signed onto a system governing 90\% of world trade, yet Article XXI had been formally invoked only once in nearly half a century.\footnote{Britannica Editors. (2026, May 7). General Agreement on Tariffs and Trade (GATT). Encyclopedia Britannica. \url{https://www.britannica.com/topic/General-Agreement-on-Tariffs-and-Trade}}

The post-1995 period saw this equilibrium unravel rapidly. China’s accession to the WTO in 2001 marked a turning point, integrating the largest emerging economy while planting seeds of the strategic rivalry that now defines global trade politics. The 2018 imposition of US Section 232 tariffs on steel and aluminium on national security grounds triggered an unprecedented wave of Article XXI invocations and WTO challenges. It exposed the inadequacy of existing rules to manage security-justified trade measures.\footnote{Reinsch, W. A., \& Caporal, J. (2025, November 17). The WTO's first ruling on national security: What does it mean for the United States? CSIS. \url{https://www.csis.org/analysis/wtos-first-ruling-national-security-what-does-it-mean-united-states}} Consequently, the Appellate Body collapsed in December 2019, removing the system’s last credible enhancement mechanism and leaving the dispute settlement architecture effectively paralysed when it was most critical. Since then, the escalation of US-China technology controls, the spread of industrial policy across major economies, and the fracturing of global supply chains along geopolitical lines have turned Article XXI from an obscure legal anomaly to the central question facing the multilateral trading system.\footnote{Maruyama, W., \& Wolff, A. Wm. (2023). Saving the WTO from the national security exception. Peterson Institute for International Economics.}
